\documentclass[11pt,a4paper]{article}

\usepackage[margin=2.45cm]{geometry}
\usepackage{amsmath,amssymb,amsthm,mathtools}
\usepackage{booktabs,longtable,array}
\usepackage{microtype}
\usepackage[hidelinks]{hyperref}
\usepackage{xcolor}
\usepackage{enumitem}

\newtheorem{theorem}{Theorem}[section]
\newtheorem{proposition}[theorem]{Proposition}
\newtheorem{lemma}[theorem]{Lemma}
\newtheorem{corollary}[theorem]{Corollary}
\newtheorem{definition}[theorem]{Definition}
\newtheorem{remark}[theorem]{Remark}
\newtheorem{problem}[theorem]{Open Problem}

\newcommand{\E}{\mathbb E}
\newcommand{\PP}{\mathbb P}
\newcommand{\cB}{\mathcal B}
\newcommand{\cK}{\mathcal K}
\newcommand{\Z}{\mathbb Z}
\newcommand{\bits}{\{0,1\}}
\newcommand{\PGF}{\operatorname{PGF}}
\newcommand{\1}{\mathbf 1}

\definecolor{provedgreen}{RGB}{20,105,55}
\definecolor{frontierred}{RGB}{145,25,25}
\newcommand{\proved}[1]{\textcolor{provedgreen}{\textbf{PROVED.} #1}}
\newcommand{\audit}[1]{\textbf{AUDITED COMPUTATION.} #1}
\newcommand{\frontier}[1]{\textcolor{frontierred}{\textbf{OPEN / ACTIVE FRONTIER.} #1}}

\title{\textbf{Who Pays the Bill?}\\[1mm]
\large Zero-Repair Checkpoint:\\
Logarithmic Divergence, Finite Binomial Viability, and Delayed Boundary Repair}
\author{Jan Mikul\'ik}
\date{22 August 2026}

\begin{document}
\maketitle

\begin{abstract}
This checkpoint isolates the present status of the exact history-indexed
zero-repair problem for the first nontrivial dyadic parity pair
\[
P=\tfrac14\delta_2+\tfrac34\delta_4,
\qquad
Q=\tfrac34\delta_3+\tfrac14\delta_5.
\]
The exact recursion is
\[
\mu_Q*B_h
=
\tfrac14\delta_2*B_{h0}
+
\tfrac34\delta_4*B_{h1},
\]
and \(\cB_n\) denotes the minimum root mean among depth-\(n\) feasible trees.
The previously established support-independent viability certificate gives
\(\E J_\varnothing\ge19/9\) for every infinite viable tree, while a direct
finite-horizon certificate already gives
\(\cB_{15}\ge39473/19683>2\).

The main quantitative update is a finite-horizon transform/Jensen argument
which yields
\[
\boxed{
\liminf_{n\to\infty}\frac{\cB_n}{\log n}\ge\frac13
}
\]
for the natural logarithm.  In particular \(\cB_n\to\infty\).
This settles the old bounded-versus-divergent viability dichotomy, but it does
\emph{not} determine the divergence exponent beyond the logarithmic lower law.

On the constructive side, small explicit binomially smoothed roots survive far
beyond their first doubling scales: one family member is analytically viable to
depth two, a second has an audited full-history tree to depth six, and a third
has an audited full-history tree to depth eleven.  More importantly, a genuine
boundary obstruction in the geometric-reservoir attack can be removed exactly.
The negative second boundary profile can be repaired by transporting a positive
packet to an arbitrarily delayed tail location \(J\), subject to one explicit
coefficient inequality.  This produces a rigorous delayed scale-switching
mechanism and an integer tangent model for the remaining renormalization
problem.

No claim of \(\cB_n=\Theta(\log n)\) or
\(\cB_n=\Theta(n^{1/3})\) is made here.
\end{abstract}

\paragraph{Status convention.}
Theorems, propositions, lemmas, and corollaries below are analytic statements.
Items labelled \textbf{AUDITED COMPUTATION} are exact finite LP feasibility
checks and are not promoted to asymptotic theorems.  Open conjectures are
explicitly marked.

\tableofcontents

% ============================================================
\section{Exact zero-repair viability problem}
% ============================================================

For each history \(h\in\bits^{\le n}\), let \(B_h\) be a probability law on
\(\Z_{\ge0}\).  Write its probability generating function as
\[
B_h(z)=\sum_{j\ge0}b_{h,j}z^j.
\]
The exact zero-repair recursion is
\[
\boxed{
\mu_Q*B_h
=
\frac14\delta_2*B_{h0}
+
\frac34\delta_4*B_{h1}.
}
\tag{ZR}
\]
Equivalently,
\[
\boxed{
B_{h0}(z)+3z^2B_{h1}(z)
=
z(3+z^2)B_h(z).
}
\tag{PGF}
\]

\begin{definition}[Finite-depth viability sets]
Let \(\cK_0=\mathcal P(\Z_{\ge0})\).  Recursively, \(B\in\cK_{n+1}\) if
there exist \(B_0,B_1\in\cK_n\) satisfying \eqref{PGF}.  Define
\[
\boxed{
\cB_n:=\inf_{B\in\cK_n}\E_BJ.
}
\]
\end{definition}

The sets are nested,
\[
\cK_{n+1}\subseteq\cK_n,
\]
so \(\cB_n\) is nondecreasing.

% ============================================================
\section{Exact lower certificates already in force}
% ============================================================

\subsection{Infinite-support viability barrier}

The closed low-tail envelope and its two-generation dual refinement give the
exact support-independent theorem
\[
\boxed{
\E_BJ\ge\frac{19}{9}
}
\]
for the root of every infinite viable zero-repair tree.
The certificate is coefficientwise on the whole support; it is not a
finite-cutoff LP artifact.

\subsection{Finite-horizon barrier above two}

A finite-depth version of the same tail hierarchy gives exact rational lower
bounds \(D_n\le\cB_n\).  In particular,
\[
\boxed{
\cB_{15}
\ge
\frac{39473}{19683}
=
2.005436\ldots
>2.
}
\]
Hence the old candidate \(\cB_\infty=2\) is already impossible at finite depth.

% ============================================================
\section{Quantitative divergence: a logarithmic theorem}
% ============================================================

Let \(m_h:=\E_{B_h}J\).  Differentiating \eqref{PGF} at \(z=1\) gives
\[
\boxed{
 m_h=\frac14m_{h0}+\frac34m_{h1}.
}
\tag{M}
\]
Thus the buffer mean is harmonic along the \(P\)-history tree.

Fix \(z\in(0,1)\), set
\[
a:=-\log z,
\qquad
y_h:=z^{m_h},
\]
and let \(H_t\) be a random history generated with branch probabilities
\(1/4,3/4\).  Convexity gives a submartingale
\[
\E[y_{H_{t+1}}\mid H_t=h]\ge y_h.
\]
For a block of length \(k\), write
\[
D_h^{(k)}
:=
\E_{P^k}y_{hw}-y_h
\ge0.
\]
Since the geometric mean of \(\sqrt{y_{hw}}\) under \(P^k\) equals
\(\sqrt{y_h}\),
\[
\boxed{
\E_{P^k}
\left(\sqrt{y_{hw}}-\sqrt{y_h}\right)^2
\le D_h^{(k)}.
}
\tag{J}
\]

Define
\[
G_P(z)=\frac{z^2+3z^4}{4},
\qquad
G_Q(z)=\frac{3z^3+z^5}{4},
\]
\[
\rho(z):=\frac{G_Q(z)}{G_P(z)}<1,
\qquad
\widetilde p_z(x):=\frac{p(x)z^x}{G_P(z)},
\]
and the exact quadratic likelihood cost
\[
\Lambda(z)
:=
\sum_x\frac{\widetilde p_z(x)^2}{p(x)}
=
\frac{G_P(z^2)}{G_P(z)^2}.
\]
Iterating \eqref{PGF} gives the exact tilted transform identity
\[
\boxed{
\E_{\widetilde P^k}B_{hw}(z)
=
\rho(z)^kB_h(z).
}
\tag{T}
\]
Using \(B_{hw}(z)\ge z^{m_{hw}}\), \eqref{J}, and Cauchy--Schwarz gives
\[
\rho(z)^kB_h(z)
\ge
 y_h-2\Lambda(z)^{k/2}\sqrt{D_h^{(k)}}.
\tag{B}
\]

Partition a depth-\(n\) path into disjoint blocks of length \(k\le n/2\).
The total submartingale increase is at most
\[
1-z^{m_\varnothing}
\le am_\varnothing.
\]
Therefore at least one block satisfies
\[
\E D_{H_s}^{(k)}
\le
\frac{2am_\varnothing k}{n}.
\]
Averaging \eqref{B} over the start of that block, using
\(B_h(z)\le1\) and \(\E y_{H_s}\ge z^{m_\varnothing}\), yields the master inequality
\[
\boxed{
\rho(z)^k
\ge
z^{m_\varnothing}
-
2\Lambda(z)^{k/2}
\sqrt{\frac{2am_\varnothing k}{n}}.
}
\tag{Master}
\]

Choose
\[
k=
\left\lceil
\frac{am_\varnothing+\log2}{-\log\rho(z)}
\right\rceil.
\]
If \(k>n/2\), then \(m_\varnothing=\Omega(n)\) and the claimed logarithmic
lower bound is immediate.  Otherwise \eqref{Master} implies
\[
\log n
\le
\Gamma(z)m_\varnothing+O_z(\log(m_\varnothing+1)),
\]
where
\[
\boxed{
\Gamma(z)
=
2a
+
\frac{a}{-\log\rho(z)}\log\Lambda(z).
}
\]
For \(z=e^{-a}\uparrow1\), exact cumulant expansion gives
\[
-\log\rho(e^{-a})
=
\frac14a^3+O(a^4),
\]
\[
\log\Lambda(e^{-a})
=
\frac34a^2+O(a^3),
\]
and hence
\[
\Gamma(e^{-a})=3+O(a).
\]

\begin{theorem}[Quantitative zero-repair divergence]
For the \(R=2\) zero-repair hierarchy,
\[
\boxed{
\liminf_{n\to\infty}
\frac{\cB_n}{\log n}
\ge
\frac13.
}
\]
Equivalently,
\[
\cB_n
\ge
\left(\frac13-o(1)\right)\log n.
\]
In particular,
\[
\boxed{\cB_n\to\infty.}
\]
\end{theorem}

\begin{remark}
The logarithm is natural.  In base two the coefficient is
\[
\frac{\log2}{3}=0.231049\ldots.
\]
The proof is finite-horizon and does not require an infinite-tree compactness
argument.
\end{remark}

% ============================================================
\section{The first two viability sets exactly}
% ============================================================

Write \(B=(b_0,b_1,\ldots)\).
Multiplying \eqref{PGF} coefficientwise by the common denominator gives
\[
3b_{u-3}+b_{u-5}
=
\ell_{u-2}+3h_{u-4},
\tag{R}
\]
with negative-index coefficients interpreted as zero.

\begin{proposition}[Exact characterization of \(\cK_1\)]
\[
\boxed{
B\in\cK_1
\iff
b_0\le\frac13.
}
\]
\end{proposition}

\begin{proof}
At \(u=3\), \eqref{R} gives \(\ell_1=3b_0\), so normalization of the low child
forces \(3b_0\le1\).  Conversely, if \(3b_0\le1\), the forced low mass can be
completed from the remaining source profile to a low submeasure of total mass
one, which is exactly the one-step Hall construction.
\end{proof}

\begin{proposition}[Exact characterization of \(\cK_2\)]
\[
\boxed{
B\in\cK_2
\iff
b_0\le\frac13,
\qquad
b_0+b_1\le\frac23.
}
\]
\end{proposition}

\begin{proof}
The low child always has \(\ell_0=0\).  From the first two nontrivial equations,
\[
\ell_1=3b_0,
\qquad
\ell_2+3h_0=3b_1.
\]
Both children must lie in \(\cK_1\), hence \(h_0\le1/3\).  Necessity follows
from
\[
3(b_0+b_1)=\ell_1+\ell_2+3h_0\le2.
\]
For sufficiency choose
\[
\ell_2=\min\{3b_1,1-3b_0\}.
\]
Then
\[
h_0=b_1-\frac13\ell_2\le\frac13
\]
precisely under \(b_0+b_1\le2/3\), and the remaining low mass can again be
completed from higher source levels.  Thus both children belong to \(\cK_1\).
\end{proof}

% ============================================================
\section{Binomial smoothing: finite-depth evidence}
% ============================================================

Let
\[
U(z):=\frac{1+z}{2},
\qquad
R_k(z):=\frac{1+2z}{3}\,U(z)^k.
\]
Then
\[
\E R_k=\frac23+\frac{k}{2}.
\]
The exact \(\cK_1,\cK_2\) characterization immediately gives
\[
R_1\in\cK_2.
\]
Moreover the finite-horizon coupled-tail inequality at depth three gives
\[
R_1\notin\cK_3.
\]
For \(R_2\), the support-independent depth-seven CDF cap
\(F(2)\le22/27\) is violated because
\[
F_{R_2}(2)=\frac56>\frac{22}{27},
\]
so
\[
R_2\notin\cK_7.
\]

The full-history feasibility LP has been rebuilt directly from \eqref{R}, with
no imported primal solver and no dual constraints.  It gives:
\[
\boxed{R_2\in\cK_6}
\]
and, in a deeper run with support cutoff \(0\le J\le16\),
\[
\boxed{R_3\in\cK_{11}}.
\]
These are positive certificates: a finite-support feasible tree is a genuine
full-history tree.  The depth-twelve status of \(R_3\) is not claimed here.

\audit{The current independent audit confirms \(R_2\in\cK_6\) and
\(R_3\in\cK_{11}\).  The negative statements above are analytic and do not rely
on finite-support LP infeasibility.}

The data suggest that binomial smoothing is useful, but they do not yet imply
any uniform doubling law such as \(R_k\in\cK_{2^k}\) for all \(k\).

% ============================================================
\section{Shift and bulk-extension lemmas}
% ============================================================

\begin{lemma}[Shift invariance]
If \(B\in\cK_n\), then \(z^sB\in\cK_n\) for every integer \(s\ge0\).
\end{lemma}

\begin{proof}
Shift every buffer law in a feasible tree by the same deterministic amount
\(s\).  Equation \eqref{PGF} is multiplied by \(z^s\) at every node.
\end{proof}

\begin{lemma}[One free bulk level]
If \(B\in\cK_n\), then
\[
\boxed{zB\in\cK_{n+1}.}
\]
\end{lemma}

\begin{proof}
The exact boundary-free split is
\[
zB\longmapsto(z^4B,B),
\]
because
\[
z^4B+3z^2B=z^2(3+z^2)B=z(3+z^2)(zB).
\]
Both children lie in \(\cK_n\) by shift invariance.
\end{proof}

% ============================================================
\section{Geometric reservoir and delayed boundary repair}
% ============================================================

Let
\[
G_r(z):=\frac{1-r}{1-rz},
\qquad 0<r<1,
\]
and define
\[
D_r(z):=(3-2z^2)G_r(z).
\]
For \(r^2\ge2/3\), \(D_r\) is a PGF and the previously known exact split is
\[
\boxed{
G_r\longmapsto(zD_r,zG_r).
}
\tag{G1}
\]
The naive next transform \((3-2z^2)D_r\) has negative coefficients.  The
following theorem repairs that defect exactly.

\begin{theorem}[Delayed boundary repair]
Set
\[
A_r:=3(1-r)(4-3r^2),
\qquad
\alpha_r:=A_r(1+r)=3(1-r^2)(4-3r^2).
\]
Fix an integer \(J\ge4\) and define
\[
C_{r,J}(z)
:=
A_rz^2+rA_rz^3-\alpha_r z^JG_r(z).
\]
Let
\[
E_{r,J}(z)
:=
(3-2z^2)^2G_r(z)+C_{r,J}(z),
\]
\[
H_{r,J}(z)
:=
zD_r(z)-\frac{C_{r,J}(z)}{3z}.
\]
Then the algebraic identity
\[
\boxed{
z(3+z^2)D_r
=
zE_{r,J}+3z^2H_{r,J}
}
\tag{DR}
\]
holds identically.

If
\[
\boxed{
\alpha_r
\le
(3r^2-2)^2r^{J-4},
}
\tag{Pos}
\]
then \(E_{r,J}\) and \(H_{r,J}\) are probability PGFs.  Hence \eqref{DR}
is a genuine positive zero-repair split.
\end{theorem}

\begin{proof}
The identity is immediate after cancellation of \(C_{r,J}\):
\[
\begin{aligned}
zE_{r,J}+3z^2H_{r,J}
&=z(3-2z^2)^2G_r+3z^3D_r\\
&=z(3-2z^2)(3+z^2)G_r.
\end{aligned}
\]
For positivity, the coefficients of \((3-2z^2)^2G_r\) at levels two and
three are exactly
\[
-A_r,
\qquad
-rA_r,
\]
so the first two terms of \(C_{r,J}\) cancel them.  For every \(n\ge J\), the
remaining coefficient of \(E_{r,J}\) is
\[
(1-r)r^{n-J}
\left[
(3r^2-2)^2r^{J-4}-\alpha_r
\right],
\]
which is nonnegative exactly under \eqref{Pos}.  Also
\(C_{r,J}(1)=0\), hence \(E_{r,J}(1)=1\).

For the high child,
\[
H_{r,J}
=(3z-2z^3)G_r
-\frac{A_r}{3}z
-\frac{rA_r}{3}z^2
+\frac{\alpha_r}{3}z^{J-1}G_r.
\]
The first two coefficients become
\[
(1-r)(3r^2-1),
\qquad
r(1-r)(3r^2-1),
\]
and all later coefficients are nonnegative for \(r^2\ge2/3\); the delayed
packet is positive.  Normalization follows from \eqref{DR} at \(z=1\).
\end{proof}

\subsection{The shortest repair}

For \(J=4\), condition \eqref{Pos} reduces exactly to
\[
9r^2-8\ge0.
\]
Thus for \(r^2\ge8/9\), one may write
\[
E_r(z)
=
9(1-r)+9r(1-r)z+(9r^2-8)z^4G_r(z),
\]
\[
H_r(z)
=
(1-r)(3r^2-1)z
+r(1-r)(3r^2-1)z^2
+(3r^4-4r^2+2)z^3G_r(z),
\]
and obtain
\[
z(3+z^2)D_r=zE_r+3z^2H_r.
\]
Both children have zero mass at level zero, hence lie in \(\cK_1\).  Therefore
\[
\boxed{D_r\in\cK_2.}
\]
Using \eqref{G1}, shift invariance, and the one-free-level lemma gives the
rigorous finite-depth corollary
\[
\boxed{G_r\in\cK_4\qquad(r^2\ge8/9).}
\]

\subsection{How far can one packet be delayed?}

Condition \eqref{Pos} permits a maximal delay satisfying
\[
r^{J-4}
\gtrsim
\frac{\alpha_r}{(3r^2-2)^2}.
\]
For \(r=1-\varepsilon\),
\[
\alpha_r=6\varepsilon+O(\varepsilon^2),
\]
so
\[
\boxed{
J_{\max}
\asymp
\frac{\log(1/\varepsilon)}{\varepsilon}.
}
\]
Since \(\E G_r=r/(1-r)\asymp\varepsilon^{-1}\), one exact local repair can
move boundary debt to distance
\[
J\asymp m\log m
\]
when the geometric mean scale is \(m\).

% ============================================================
\section{Integer tangent model}
% ============================================================

Divide the geometric profiles by \(1-r\) and let \(r\uparrow1\).  Coefficientwise,
\[
\frac{G_r(z)}{1-r}
\longrightarrow
\frac1{1-z},
\]
whose coefficient sequence is
\[
(1,1,1,1,\ldots).
\]
After one boundary transform,
\[
\frac{D_r(z)}{1-r}
\longrightarrow
\frac{3-2z^2}{1-z},
\]
with coefficients
\[
(3,3,1,1,1,\ldots).
\]
After the naive second transform,
\[
\frac{(3-2z^2)^2G_r(z)}{1-r}
\longrightarrow
\frac{(3-2z^2)^2}{1-z},
\]
with coefficients
\[
\boxed{
(9,9,-3,-3,1,1,1,\ldots).
}
\]
Thus the small-\((1-r)\) boundary problem has a discrete tangent description:
two negative packets of size three must be routed to later positive tail scales.
The delayed correction above performs exactly such a routing before the tangent
limit is taken.

% ============================================================
\section{What is and is not settled}
% ============================================================

The lower side is now rigorous:
\[
\boxed{
\cB_n
\ge
\left(\frac13-o(1)\right)\log n,
\qquad
\cB_n\to\infty.
}
\]
This rules out every bounded exact zero-repair hierarchy.

The constructive side remains unresolved.  In particular, neither
\[
\cB_n=O(\log n)
\]
nor
\[
\cB_n=O(n^{1/3})
\]
has been proved for the exact zero-repair problem.  The atomic B-spline
\(O(n^{1/3})\) construction belongs to the broader causal-bill problem with
positive residual repair and must not be imported as a zero-repair upper bound.

The exact delayed repair theorem shows that the boundary defect can be moved
between scales rather than paid immediately.  The remaining issue is whether
these packets can be routed through all histories with an accumulated root mean
of order \(\log n\), a power law, or another scale.

\begin{problem}[Doubling increment]
Determine the asymptotic behavior of
\[
\boxed{
\cB_{2n}-\cB_n.
}
\]
A uniform bound would imply \(\cB_n=O(\log n)\); divergence of the increment
would rule out that route.
\end{problem}

\begin{problem}[Boundary-debt renormalization]
Construct a positive multiscale state variable which records the delayed
packets created by \eqref{DR}, and determine the Perron/renormalization factor
of its tangent dynamics.  The integer profile
\[
(9,9,-3,-3,1,1,\ldots)
\]
is the simplest local model of the defect to be transported.
\end{problem}

% ============================================================
\section{Current ledger}
% ============================================================

\small
\begin{longtable}{p{0.72\textwidth}|p{0.20\textwidth}}
\textbf{Statement} & \textbf{Status}\\
\hline
\endfirsthead
\textbf{Statement} & \textbf{Status}\\
\hline
\endhead
Exact zero-repair PGF recursion & proved\\
Monotonicity of \(\cB_n\) & proved\\
Infinite-tree lower bound \(\E J\ge19/9\) & proved\\
Finite-depth bound \(\cB_{15}\ge39473/19683>2\) & proved\\
\(\liminf \cB_n/\log n\ge1/3\) & \textbf{proved}\\
\(\cB_n\to\infty\) & \textbf{proved}\\
Exact characterization of \(\cK_1\) and \(\cK_2\) & proved\\
\(R_1\in\cK_2\), \(R_1\notin\cK_3\) & proved\\
\(R_2\in\cK_6\) & audited computation\\
\(R_2\notin\cK_7\) & proved by CDF cap\\
\(R_3\in\cK_{11}\) & audited computation\\
Delayed boundary-repair identity \eqref{DR} & proved\\
Positivity condition \eqref{Pos} & proved\\
\(D_r\in\cK_2\), \(G_r\in\cK_4\) for \(r^2\ge8/9\) & proved\\
Uniform logarithmic upper \(\cB_n=O(\log n)\) & \textcolor{frontierred}{open}\\
Cubic exact-zero-repair law \(\cB_n=\Theta(n^{1/3})\) & \textcolor{frontierred}{open}\\
Full asymptotic order of \(\cB_n\) & \textcolor{frontierred}{open}
\end{longtable}
\normalsize

\section*{Bottom line}
The old finite-mean viability branch is closed: exact zero repair has a
quantitative logarithmic divergence.  The present frontier is therefore no
longer existence versus nonexistence.  It is the growth law.  The strongest new
constructive information is that a genuine second-order boundary negativity is
not terminal: it can be repaired exactly by exporting a positive packet to a
later geometric scale.  The next proof target is the induced multiscale packet
dynamics, not another one-scale stationary ansatz.

\end{document}
